\documentclass[11pt]{article}
\usepackage[margin=1in]{geometry}
\usepackage{amsmath,amssymb,amsthm}
\usepackage{graphicx}
\usepackage{booktabs}
\usepackage{xcolor}
\usepackage{listings}
\usepackage[numbers]{natbib}
\usepackage{hyperref}

\newtheorem{theorem}{Theorem}
\newtheorem{lemma}[theorem]{Lemma}
\theoremstyle{definition}
\newtheorem{definition}[theorem]{Definition}
\theoremstyle{remark}
\newtheorem*{remark}{Remark}

\emergencystretch=2em

\lstdefinelanguage{lean}{
  morekeywords={def,theorem,rcases,haveI,subst,nlinarith,linarith,have,refine,intro,induction,with,simp,only,rw,unfold,exact,example,namespace,end,open,fun,Prop,by,noncomputable,section,import},
  sensitive=true,
  morecomment=[l]{--},
  morecomment=[s]{/-}{-/},
}
\lstset{
  language=lean,
  basicstyle=\ttfamily\footnotesize,
  keywordstyle=\color{blue!60!black}\bfseries,
  commentstyle=\color{gray},
  columns=fullflexible,
  keepspaces=true,
  breaklines=true,
  frame=single,
  extendedchars=true,
  literate=
    {ℝ}{{$\mathbb{R}$}}1
    {ℕ}{{$\mathbb{N}$}}1
    {→}{{$\to$}}1
    {↔}{{$\leftrightarrow$}}1
    {∀}{{$\forall$}}1
    {∃}{{$\exists$}}1
    {∈}{{$\in$}}1
    {∑}{{$\sum$}}1
    {∏}{{$\prod$}}1
    {∧}{{$\wedge$}}1
    {≤}{{$\le$}}1
    {≠}{{$\ne$}}1
    {⨆}{{$\bigsqcup$}}1
    {∉}{{$\notin$}}1
    {←}{{$\leftarrow$}}1
    {ρ}{{$\rho$}}1
    {⟨}{{$\langle$}}1
    {⟩}{{$\rangle$}}1
    {·}{{$\cdot$}}1
    {𝓝}{{$\mathcal{N}$}}1
    {₀}{{$_0$}}1
    {⊢}{{$\vdash$}}1
}

\title{Cascaded Recovery Bounds:\\ A Lean Proof that a Cascade of Imperfect Stages Beats Each of Its Stages}
\author{Justin Hart, Viridis LLC}
\date{16 November 2026}

\begin{document}
\maketitle

\begin{abstract}
A material stream passes through $n$ recovery stages in series. Stage $i$ recovers a fraction
$r_i\in[0,1]$ of what reaches it. The share of the original stream that must still be met by
virgin material, the \emph{virgin share}, is $\prod_i(1-r_i)$. We give a Lean~4 proof, against
Mathlib, of the sealed statement \texttt{P3.CascadedRecovery}. The statement has four parts.
The virgin share lies in $[0,1]$. It is at least $(1-m)^n$ for every uniform upper bound $m\le1$ on
the stage rates, and in particular at least $(1-\max_i r_i)^n$. It is never larger than the
single-stage share $1-r_j$ of any stage $j$. It is strictly smaller than $1-r_j$ whenever
$r_j<1$ and some other stage recovers a positive fraction. The proof is about seventy lines and
rests on one idea: split stage $j$ off the product and bound the rest by $1$. These results are
properties of the series-cascade model. They make no claim about any real stage efficiency,
material or recovery technology.
\end{abstract}

\section{Introduction}
Recovery processes are rarely perfect. A sorting line misses some material, a re-processing step
rejects some of what it receives, and a second or third pass is often added to catch what the
first one missed. The circular-economy literature treats closing material loops as the core of
the concept \citep{geissdoerfer2017}. Industrial symbiosis, where one firm's residue becomes
another's input, extends this to chains of firms \citep{chertow2000}. In both settings a natural
question for a decision kernel is: \emph{if several imperfect recovery stages act in series, how
much virgin material is still needed, and is the cascade always at least as good as its best
single stage?}

In the series model the answer is a product, and the facts are elementary. Yet the edge cases
matter for a kernel that applies them automatically. What happens with zero stages? With a stage
that recovers everything? With a stage that recovers nothing? When is the improvement strict?
Project P3 of Viridis Research proposed the statement \texttt{P3.CascadedRecovery} in session
P3-T001 (5~November 2026). The program owner sealed it under decision D-012. This note gives a
Lean proof that the sealed statement holds exactly as written.

\paragraph{Contributions.}
\begin{itemize}
\item A Lean~4 proof of \texttt{ViridisStatements.P3.CascadedRecovery.Statement}
(Section~\ref{sec:proof}). It depends only on the standard axioms \texttt{propext},
\texttt{Classical.choice} and \texttt{Quot.sound}.
\item A short human-readable proof (Section~\ref{sec:math}) that shows which hypothesis each
part uses, including the role of Mathlib's convention for the supremum over an empty index set.
\item An account of what the result settles and what it does not (Section~\ref{sec:limits}).
\end{itemize}

\section{Background and related work}
\paragraph{Circular economy and cascades.}
\citet{geissdoerfer2017} review how the circular economy relates to sustainability. Its
operational core is substituting recovered material for virgin extraction. \citet{chertow2000}
defined industrial symbiosis as the exchange of materials, energy, water and by-products among
firms. Such exchanges often form chains, where the residue of one stage feeds the next. The
previous P3 session (P3-T002, 15~November 2026) gave a Lean proof of the single-loop mass balance
\texttt{P3.LoopMassBalance}, where virgin input is $D(1-r)$ for one recovery fraction $r$. The
cascade model here generalises the factor $1-r$ to a product over stages.

\paragraph{The mathematics.}
That pass-through fractions of stages in series multiply, and that a product of numbers in
$[0,1]$ is bounded by each factor, are textbook facts. The same structure appears in the
reliability of series systems and in serial separation efficiencies. We claim no mathematical
novelty. The contribution is a machine-checkable proof of exactly the sealed statement,
including its strictness condition and its degenerate cases.

\paragraph{Formal tools.}
The proof is written in Lean~4 \citep{moura2021lean4} (toolchain v4.28.0) against Mathlib
\citep{mathlib2020} (commit \texttt{8f9d9cff}).

\section{The model and the statement}
\label{sec:model}
\begin{definition}[Cascade model]
Let $n\in\mathbb{N}$ and $r:\{0,\dots,n-1\}\to\mathbb{R}$. The \emph{virgin share} of the cascade
is
\[
\mathrm{VS}(r)=\prod_{i<n}(1-r_i).
\]
\end{definition}
The reading is as follows. One unit of material enters stage $0$. Stage $i$ recovers the fraction
$r_i$ of what reaches it and passes the fraction $1-r_i$ on unrecovered. After all $n$ stages the
unrecovered fraction is $\mathrm{VS}(r)$, and this is the share of demand that virgin material
must cover.

\begin{theorem}[\texttt{P3.CascadedRecovery}]\label{thm:main}
Let $0\le r_i\le1$ for all $i$. Then:
\begin{enumerate}
\item $0\le\mathrm{VS}(r)\le1$;
\item for every $m$ with $r_i\le m$ for all $i$ and $m\le1$: $(1-m)^n\le\mathrm{VS}(r)$;
\item $(1-\sup_i r_i)^n\le\mathrm{VS}(r)$, where $\sup$ is Mathlib's $\bigsqcup_i r_i$ in the
conditionally complete lattice $\mathbb{R}$;
\item for every stage $j$: $\mathrm{VS}(r)\le 1-r_j$;
\item for all stages $j\ne k$ with $0<r_k$ and $r_j<1$: $\mathrm{VS}(r)<1-r_j$.
\end{enumerate}
\end{theorem}

The sealed Lean text is reproduced below.

\begin{lstlisting}[caption={Definition and \texttt{Statement} from the sealed module \texttt{ViridisStatements.P3.CascadedRecovery} (docstrings omitted).}]
def virginShare {n : ℕ} (r : Fin n → ℝ) : ℝ := ∏ i, (1 - r i)

def Statement : Prop :=
  ∀ (n : ℕ) (r : Fin n → ℝ), (∀ i, 0 ≤ r i) → (∀ i, r i ≤ 1) →
    (0 ≤ virginShare r ∧ virginShare r ≤ 1) ∧
    (∀ m : ℝ, (∀ i, r i ≤ m) → m ≤ 1 → (1 - m) ^ n ≤ virginShare r) ∧
    (1 - ⨆ i, r i) ^ n ≤ virginShare r ∧
    (∀ j, virginShare r ≤ 1 - r j) ∧
    (∀ j k, k ≠ j → 0 < r k → r j < 1 → virginShare r < 1 - r j)
\end{lstlisting}

\section{Mathematical proof}
\label{sec:math}
Write $f_i=1-r_i$. The hypotheses give $0\le f_i\le1$ for every $i$.

\begin{lemma}[Splitting off a stage]\label{lem:split}
For every stage $j$, $\mathrm{VS}(r)=f_j\cdot P_j$, where $P_j=\prod_{i\ne j}f_i\in[0,1]$.
\end{lemma}
\begin{proof}
Commutativity of the product gives the identity. A product of numbers in $[0,1]$ lies in
$[0,1]$.
\end{proof}

\begin{proof}[Proof of Theorem~\ref{thm:main}]
(1) Every factor lies in $[0,1]$, so the product does too.

(2) From $r_i\le m\le1$ we get $0\le1-m\le f_i$. The product of $n$ copies of $1-m$ is
$(1-m)^n$. Products of non-negative numbers are monotone in each factor, so
$(1-m)^n\le\prod_i f_i$. This part does not need $r_i\ge0$.

(3) If $n=0$, the empty product is $1$ and $(1-s)^0=1$ for any $s$, so the claim holds with
equality. Mathlib's value of $\bigsqcup$ over the empty type, which is $0$, never matters. If
$n>0$, the index set is finite and non-empty, so $s=\sup_i r_i$ is a true maximum. Then $r_i\le
s$ for all $i$ and $s\le1$, and part~(2) with $m=s$ applies.

(4) By Lemma~\ref{lem:split}, $\mathrm{VS}(r)=f_jP_j\le f_j$, because $f_j\ge0$ and $P_j\le1$.

(5) Since $k\ne j$, split $k$ off $P_j$: $P_j=f_k\cdot Q$ with $Q=\prod_{i\ne j,k}f_i\in[0,1]$.
Then $P_j\le f_k=1-r_k<1$. Since $r_j<1$, $f_j>0$, so $\mathrm{VS}(r)=f_jP_j<f_j$.
\end{proof}

\begin{remark}
The sealed module's negative controls show that every hypothesis of part~(5) is needed. If
every other stage recovers nothing, the cascade equals stage $j$ alone. If $r_j=1$, both sides
are $0$. Both bounds $0\le r_i$ and $r_i\le1$ are needed too. A ``stage'' with $r=2$ makes the
share negative. A ``stage'' with $r=-1$, one that adds demand, pushes the share above another
stage's share.
\end{remark}

\paragraph{Interpretation.} Parts~(4) and~(5) say that adding stages never hurts and that it
helps strictly as soon as another stage does any work. So a cascade of imperfect stages beats
any one of its stages, including the best one. Parts~(2) and~(3) bound how much the cascade can
achieve. With every stage at most as good as the best one, rate $m$, the virgin share cannot fall
below what $n$ copies of that best stage would leave, which is $(1-m)^n$. For example, two stages
at $1/2$ give share $1/4$: equal to the lower bound and strictly below the single-stage share
$1/2$. The stages $1/2,1/4,0$ give $3/8$, between $(1/2)^3=1/8$ and $1/2$.

\section{The Lean development}
\label{sec:proof}
The candidate file \texttt{proofs/P3.CascadedRecovery.lean} follows Section~\ref{sec:math}.
\texttt{subprod\_le\_one} and \texttt{subprod\_nonneg} bound any sub-product, via
\texttt{Finset.prod\_le\_one} and \texttt{Finset.prod\_nonneg}. \texttt{split} is
Lemma~\ref{lem:split}, which is \texttt{Finset.mul\_prod\_erase}. \texttt{lower\_bound} is
part~(2). It rewrites $(1-m)^n$ as a constant product with \texttt{Finset.prod\_const} and
\texttt{Fintype.card\_fin}, then applies \texttt{Finset.prod\_le\_prod}. In \texttt{main}, part~(3)
splits on $n=0$. For $n>0$ it uses \texttt{le\_ciSup} with boundedness from
\texttt{Set.finite\_range}, and \texttt{ciSup\_le} with non-emptiness. Parts~(4) and~(5) are closed
by \texttt{nlinarith} after the split, with stage $k$ split off a second time in part~(5).

\begin{lstlisting}[caption={The complete candidate proof \texttt{proofs/P3.CascadedRecovery.lean}.}]
import Mathlib
import ViridisStatements.P3.CascadedRecovery

namespace Candidate

open ViridisStatements.P3.CascadedRecovery

/-- Every factor `1 - r i` lies in `[0, 1]`. -/
theorem factor_nonneg {n : ℕ} (r : Fin n → ℝ) (h1 : ∀ i, r i ≤ 1) (i : Fin n) :
    0 ≤ 1 - r i := by linarith [h1 i]

theorem factor_le_one {n : ℕ} (r : Fin n → ℝ) (h0 : ∀ i, 0 ≤ r i) (i : Fin n) :
    1 - r i ≤ 1 := by linarith [h0 i]

/-- Any sub-product of the factors lies in `[0, 1]`. -/
theorem subprod_le_one {n : ℕ} (r : Fin n → ℝ) (h0 : ∀ i, 0 ≤ r i) (h1 : ∀ i, r i ≤ 1)
    (s : Finset (Fin n)) : ∏ i ∈ s, (1 - r i) ≤ 1 :=
  Finset.prod_le_one (fun i _ => factor_nonneg r h1 i) (fun i _ => factor_le_one r h0 i)

theorem subprod_nonneg {n : ℕ} (r : Fin n → ℝ) (h1 : ∀ i, r i ≤ 1)
    (s : Finset (Fin n)) : 0 ≤ ∏ i ∈ s, (1 - r i) :=
  Finset.prod_nonneg (fun i _ => factor_nonneg r h1 i)

/-- Splitting off stage `j`: `virginShare r = (1 - r j) · ∏_{i ≠ j} (1 - r i)`. -/
theorem split {n : ℕ} (r : Fin n → ℝ) (j : Fin n) :
    virginShare r = (1 - r j) * ∏ i ∈ Finset.univ.erase j, (1 - r i) :=
  (Finset.mul_prod_erase Finset.univ (fun i => 1 - r i) (Finset.mem_univ j)).symm

/-- Uniform lower bound `(1 - m)^n ≤ virginShare r`. -/
theorem lower_bound {n : ℕ} (r : Fin n → ℝ) (m : ℝ) (hm : ∀ i, r i ≤ m) (hm1 : m ≤ 1) :
    (1 - m) ^ n ≤ virginShare r := by
  have hc : (1 - m) ^ n = ∏ _i : Fin n, (1 - m) := by
    rw [Finset.prod_const, Finset.card_univ, Fintype.card_fin]
  rw [hc, virginShare]
  exact Finset.prod_le_prod (fun _ _ => by linarith) (fun i _ => by linarith [hm i])

theorem main : ViridisStatements.P3.CascadedRecovery.Statement := by
  intro n r h0 h1
  refine ⟨⟨subprod_nonneg r h1 _, subprod_le_one r h0 h1 _⟩,
    fun m hm hm1 => lower_bound r m hm hm1, ?_, ?_, ?_⟩
  · -- the supremum bound
    rcases Nat.eq_zero_or_pos n with hn | hn
    · subst hn
      simp [virginShare]
    · haveI : Nonempty (Fin n) := ⟨⟨0, hn⟩⟩
      have hbdd : BddAbove (Set.range r) := (Set.finite_range r).bddAbove
      exact lower_bound r _ (fun i => le_ciSup hbdd i) (ciSup_le h1)
  · intro j
    rw [split r j]
    have hP := subprod_le_one r h0 h1 (Finset.univ.erase j)
    have hj := factor_nonneg r h1 j
    nlinarith
  · intro j k hkj hk hj
    rw [split r j]
    have hkmem : k ∈ Finset.univ.erase j := Finset.mem_erase.mpr ⟨hkj, Finset.mem_univ k⟩
    have hsplit := Finset.mul_prod_erase (Finset.univ.erase j) (fun i => 1 - r i) hkmem
    have hQ := subprod_le_one r h0 h1 ((Finset.univ.erase j).erase k)
    have hQ0 := subprod_nonneg r h1 ((Finset.univ.erase j).erase k)
    have hP : ∏ i ∈ Finset.univ.erase j, (1 - r i) < 1 := by
      rw [← hsplit]
      have hk1 := factor_nonneg r h1 k
      nlinarith
    have hj' : 0 < 1 - r j := by linarith
    nlinarith

end Candidate
\end{lstlisting}

\section{Results}
\label{sec:results}
We ran the program's checker on the candidate file. It runs the same stages as the verifier:
denylist, compile, statement match, axioms and kernel replay. The output was
\texttt{PASSED (preview)}, and the reported axioms were \texttt{propext},
\texttt{Classical.choice} and \texttt{Quot.sound}. The proof worked on the first attempt,
following the sketch recorded in the P3-T002 run's next steps. No automated prover was used. The
certification page that precedes this paper gives the verifier's own judgement.

\section{What this settles and what it does not}
\label{sec:limits}
\paragraph{Settled.} Assuming the verifier accepts the file, the bounds of the series-cascade
model hold exactly as sealed, for every number of stages including zero and for every rate
vector in $[0,1]^n$. A kernel built on this definition may rely on them.

\paragraph{Not settled.}
\begin{itemize}
\item \textbf{No empirical content.} Nothing here estimates or constrains any real stage
efficiency, material, or recovery technology. Stage rates are inputs. The theorem only says what
follows once they are given.
\item \textbf{Independence of stages.} The model assumes each stage recovers a fixed fraction of
what reaches it, whatever earlier stages did. Real stages may be correlated. For example, the
material a first sorter misses may be exactly the material a second sorter also misses, and then
the real cascade does worse than the product. The theorem says nothing about such cases.
\item \textbf{No costs, energy or quality.} Extra stages cost energy, money and possibly
material quality. The statement ``the cascade beats each stage'' concerns virgin share only. It
is not a claim that adding stages is worthwhile. Recovery work is the subject of the separate
statement \texttt{P3.RecoveryWorkFloor}.
\item \textbf{Static model.} There are no stocks or time lags. The cascade is a one-pass map;
coupling it to the loop dynamics of \texttt{P3.LoopMassBalance} (effective recovery
$r=1-\mathrm{VS}(r)$) is left for later work and is not part of any sealed statement.
\end{itemize}

\section{Reproducibility}
\sloppy
The proof file \texttt{proofs/P3.CascadedRecovery.lean} imports only \texttt{Mathlib} and the
sealed module \texttt{ViridisStatements.P3.CascadedRecovery}. It compiles under toolchain
\texttt{leanprover/lean4:v4.28.0} and the Mathlib commit below. The checker invocation is also
shown below.
\begin{lstlisting}[language={},frame=none]
sealed module sha256:
  a4e3f3c269b81a1d5dacde1535f0112654d5d4c5203e9a01a444bc4edcfbb115
Mathlib commit:
  8f9d9cff6bd728b17a24e163c9402775d9e6a365
viridis-lean-check proofs/P3.CascadedRecovery.lean P3.CascadedRecovery
\end{lstlisting}

\section{Changelog}
\begin{itemize}
\item 2026-11-05 (P3-T001): first version. Proposed P3.LoopMassBalance, P3.CascadedRecovery,
P3.RecoveryWorkFloor and P3.CyclicSymbiosis as statement modules, with controls, probe results
and rationale.
\item 2026-11-15 (P3-T002): P3.LoopMassBalance sealed (D-011). Narrowed the paper to a note on
this statement with a full Lean proof and a human-readable proof.
\item 2026-11-16 (P3-T003): P3.CascadedRecovery sealed (D-012). This version is a note on the
cascade statement with a full Lean proof. The single-loop result from P3-T002 is now summarised
in the background section. Its proof stays in the 2026-11-15 run package.
\end{itemize}

\bibliographystyle{plainnat}
\bibliography{refs}
\end{document}
